\section{Draft Length Sensitivity and Intra-Draft Refresh}
\label{appendix:draft_length}

As draft length increases, the sparse context selected at the last verification boundary becomes increasingly stale, potentially degrading acceptance rate. The intra-draft refresh mechanism (Section~\ref{section:refresh_layer}) is designed to counteract this by updating the sparse context at every draft step. To quantify its effect, we compare two variants on LB[16k-18k] (Qwen3-8B, $B{=}24$, sparsity rate $\rho = 7\%$, $L = 8$), sweeping the fixed draft length $\gamma$ from 1 to 10: one with the refresh layer enabled (as in ASPIRE), and one without, where the sparse context is updated only at verification boundaries.

Figure~\ref{fig:refresh_ablations} shows the results. At $\gamma{=}1$ both variants achieve near-identical acceptance (${\sim}$95\%), since the context has not yet gone stale. As draft length grows, the gap widens steadily: without refresh, acceptance drops to 60.3\% at $\gamma{=}10$, whereas with refresh it remains at 76.6\%, a 16.4~pp gap. This confirms that intra-draft refresh substantially mitigates context staleness, enabling longer and more effective speculation.
